%========================================================
% CHAPTER 2: LITERATURE REVIEW
%========================================================
\chapter{Literature Review}
\label{chap:literature}

\section{CAPM and Time-Varying Risk}

The Capital Asset Pricing Model, independently developed by \citet{sharpe1964capital} and \citet{lintner1965valuation}, establishes that in equilibrium, the expected excess return of an asset is proportional to its covariance with the market portfolio. Despite well-documented empirical anomalies, the CAPM remains the dominant framework for connecting risk and return in both research and practice. The Fama-French three-factor model \citep{fama1993common} and its five-factor extension \citep{fama2015five} added size (SMB), value (HML), profitability (RMW), and investment (CMA) factors to better describe the cross-section of returns. These extensions treat beta as fixed over time, an assumption that has been questioned by a growing literature on time-varying risk. \citet{harvey1989time} and \citet{ferson1996measuring} demonstrate that conditional betas fluctuate with the business cycle, while \citet{ang2007capm} show that static CAPM fails precisely during periods of high market stress. Any firm attribute that systematically alters a firm's risk profile --- including active ESG transition --- should, in principle, affect its time-varying beta.

\section{ESG and Equity Returns: The Static Impasse}

The ESG-return literature is characterised by heterogeneity. Early work by \citet{renneboog2008socially} and \citet{hong2009price} suggested that socially responsible investing imposes a financial cost, as ``sin stocks'' excluded from ESG portfolios earn superior returns. More recent meta-analytic evidence \citep{friede2015esg} finds that the majority of studies report a non-negative ESG-return relationship. \citet{edmans2011does} documents a positive long-run return to high employee satisfaction firms. \citet{khan2016corporate} find that only \emph{material} ESG characteristics predict returns, while immaterial ESG scores do not --- an early indication that the blanket aggregate score conflates relevant and irrelevant information.

The aggregate confusion is further compounded by ESG rating disagreement. \citet{berg2022aggregate} document substantial divergence across major ESG rating providers for the same firms, which mechanically reduces the statistical power of any cross-sectional ESG-return test. When ratings disagree, level-based cross-sectional tests are noisy by construction.

\citet{pastor2021sustainable} provide a theoretical reconciliation: ESG-rated firms earn lower expected returns in equilibrium because ESG-motivated investors bid up their prices, yet they show higher realised returns during periods of increasing investor ESG taste due to re-rating demand. This framework implies that the return to ESG is inherently dynamic --- driven by changes in investor preferences and ESG score revisions --- rather than by a static level premium. This thesis builds on precisely this insight.

\section{ESG Dynamics, Momentum, and Materiality}

The treatment of ESG as a dynamic variable is a recent development. \citet{giese2019foundations} document that ESG score revisions by major rating agencies are associated with abnormal returns around the event window, consistent with a price discovery channel. \citet{serafeim2020public} show that firms experiencing positive ESG momentum subsequently outperform those with flat or declining scores. \citet{lioui2021esg} argue that ESG momentum proxies for an omitted information channel: rating agencies incorporate information about future operational risk that is only gradually impounded in prices.

The ESG momentum channel connects naturally to the broader momentum literature. \citet{jegadeesh1993returns} document that past-winner stocks outperform past-losers over intermediate horizons. The delayed information diffusion mechanism of \citet{hong1999unified} --- investor underreaction to new information --- applies directly to ESG: if ESG changes convey information that institutional funds incorporate only gradually due to mandate constraints and rebalancing costs, lagged ESG changes should predict future returns.

The materiality dimension of ESG is formalised in the Sustainability Accounting Standards Board (SASB) framework, which identifies industry-specific ESG issues as financially material or immaterial based on their expected impact on enterprise value. \citet{khan2016corporate} provide the key empirical validation: only SASB-material ESG characteristics predict future returns, while immaterial ones carry no signal. This finding motivates the sector-level heterogeneity analysis in Chapter~\ref{chap:results}: ESG improvements in capital-intensive sectors such as Materials and Industrials represent genuine mitigation of regulatory and environmental risk and should be priced; ESG improvements in Information Technology may be perceived as immaterial CSR activity and should not.

Finally, the ``transition risk'' concept \citep{christophersen2022transition} describes the financial risks borne by firms during the process of improving their ESG profile: capital expenditure on new infrastructure, operational restructuring, and reputational uncertainty during the transition. This literature motivates our Beta Amplification Paradox hypothesis: firms actively improving ESG face elevated operational uncertainty precisely because they are in transition, which may temporarily increase their systematic risk loading rather than reduce it.